\documentclass[11pt]{article}
\usepackage[T1]{fontenc}
\usepackage[utf8]{inputenc}
\usepackage[margin=1in]{geometry}
\usepackage{amsmath,amssymb,amsthm}
\usepackage{microtype}
\usepackage[hidelinks]{hyperref}
\newcommand{\E}{\mathbb E}
\newcommand{\Pb}{\mathbb P}
\newcommand{\Ind}[1]{\mathbf 1\{#1\}}
\newtheorem{proposition}{Proposition}
\newtheorem{lemma}{Lemma}
\title{Persistence of High-Quality Product Use under Forgetting}
\author{Standalone theoretical note}
\date{}
\begin{document}
\maketitle

This note presents the proofs of the discussed proposition about the frequent use of the high-quality product.


\paragraph{Model and notation.}
Fix $\rho\in(0,1)$, $\gamma\in(0,1)$, $p_0\in(0,1)$,
$0<p_1<p_2<1$, and $0\leq c_1<c_2<R$.
Write $\Delta c=c_2-c_1$ and $\Delta p=p_2-p_1$.
As in the manuscript, the normalized state space and outcome transitions are
\begin{align*}
\mathcal Z&=\{(x_s,x_f):x_s\geq0,\ x_f\geq0,\ x_s+x_f\leq1\},\\
T_\rho^s(x_s,x_f)&=(\rho x_s+1-\rho,\rho x_f),\\
T_\rho^f(x_s,x_f)&=(\rho x_s,\rho x_f+1-\rho).
\end{align*}

\begin{lemma}[Continuity of $V_\rho$]For every $\rho\in[0,1)$, the Bellman equation has a unique bounded continuous solution $V_\rho$ on $\mathcal Z$.
\end{lemma}

\begin{proof}
Let $C(\mathcal Z)$ be the space of continuous real-valued functions on
$\mathcal Z$, with norm $\|v\|_\infty=\sup_{z\in\mathcal Z}|v(z)|$.
Continuity of demand and the transition maps implies that the Bellman
operator $\mathcal T_\rho$ maps $C(\mathcal Z)$ into itself. Since
\begin{equation*}
0\leq\frac{D_\rho(z)}{1-\gamma[1-D_\rho(z)]}\leq1,
\end{equation*}
we have
\begin{equation*}
\|\mathcal T_\rho v-\mathcal T_\rho w\|_\infty
\leq\gamma\|v-w\|_\infty.
\end{equation*}
Thus $\mathcal T_\rho$ is a contraction. As $\mathcal Z$ is compact,
$C(\mathcal Z)$ is complete, and the Banach fixed-point theorem gives
a unique bounded continuous solution $V_\rho$.
\end{proof}

\paragraph{The boundary set.}
For $0<\rho<1$, define
\begin{equation}
\mathcal K_\rho=
\left\{
\left((1-\rho)\sum_{j=0}^{\infty}\rho^j b_j,
      (1-\rho)\sum_{j=0}^{\infty}\rho^j(1-b_j)\right)
:\ b_j\in\{0,1\}\text{ for every }j\geq0
\right\}.
\end{equation}
Here $b_j=1$ represents a success and $b_j=0$ a failure, with $j=0$
representing the most recent outcome. Since
\begin{equation*}
(1-\rho)\sum_{j=0}^{\infty}\rho^j=1,
\end{equation*}
every point of $\mathcal K_\rho$ lies on the boundary segment
$x_s+x_f=1$. Equivalently,
\begin{equation}
\mathcal K_\rho=\{(x,1-x):x\in C_\rho\},
\qquad
C_\rho=\left\{(1-\rho)\sum_{j=0}^{\infty}\rho^j b_j:
b_j\in\{0,1\}\text{ for every }j\geq0\right\}.
\end{equation}
We call $\mathcal K_\rho$ the limiting reachable set: its points can be
approached arbitrarily closely by states generated from $(0,0)$ through
increasingly long histories of successes and failures. It satisfies
\begin{equation}
\mathcal K_\rho=T_\rho^s(\mathcal K_\rho)\cup T_\rho^f(\mathcal K_\rho).
\end{equation}
If $\rho\geq1/2$, then $C_\rho=[0,1]$, so $\mathcal K_\rho$ is the
entire boundary segment. If $0<\rho<1/2$, then $C_\rho$ is a Cantor set,
so $\mathcal K_\rho$ is only a subset of that segment.

Thus $\mathcal K_\rho$ describes limiting reachable states, rather than
states reached after finitely many observations. In particular, from
$(0,0)$ and for $0<\rho<1$, the coordinate sum after $n$ observations is
$1-\rho^n<1$. Such states approach the boundary without reaching it in
finite time. We next verify compactness of the limiting set.

\begin{lemma}[Compactness of $\mathcal{K}_\rho$] For $\rho \in (0,1)$ we have that the set $\mathcal{K}_\rho$ is compact. 
\end{lemma}

\begin{proof}
On the boundary, write $(x_s,x_f)=(x,1-x)$. Start with all possible
first coordinates, $K_0=[0,1]$, then 
\begin{equation*}
K_{n+1}=f_0(K_n)\cup f_1(K_n),
\qquad f_0(x)=\rho x,\quad f_1(x)=\rho x+(1-\rho).
\end{equation*}
Since $K_0$ is compact and both maps are continuous, induction shows
that every $K_n$ is compact. Moreover, $K_1\subseteq K_0$; applying
both maps to each inclusion gives
\begin{equation*}
K_{n+1}\subseteq K_n\qquad\text{for every }n\geq0.
\end{equation*}
We now prove the equality
\begin{equation*}
C_\rho=\bigcap_{n=0}^{\infty}K_n.
\end{equation*}

\medskip
\noindent\textit{First inclusion: $C_\rho\subseteq\bigcap_{n\geq0}K_n$.}
Let $x\in C_\rho$. By definition, there are digits $b_j\in\{0,1\}$
such that $x=(1-\rho)\sum_{j=0}^{\infty}\rho^j b_j$.
For any $n\geq1$, separate the first $n$ terms:
\begin{equation*}
x=(1-\rho)\sum_{j=0}^{n-1}\rho^j b_j+\rho^n y_n,
\qquad
y_n=(1-\rho)\sum_{k=0}^{\infty}\rho^k b_{n+k}\in[0,1].
\end{equation*}
This is precisely the result of applying $n$ transitions to $y_n$:
\begin{equation*}
x=f_{b_0}\circ f_{b_1}\circ\cdots\circ f_{b_{n-1}}(y_n).
\end{equation*}
Since $y_n\in K_0$, it follows that $x\in K_n$. This holds for every
$n$, and $x\in K_0$ as well, so $x\in\bigcap_{n\geq0}K_n$.

\medskip
\noindent\textit{Second inclusion: $\bigcap_{n\geq0}K_n\subseteq C_\rho$.}
Let $x\in\bigcap_{n\geq0}K_n$. Membership in each $K_n$ means that
we can represent this same fixed number using $n$ digits and a remainder.
Choose one such representation for each $n$. The first three have the form
\begin{align*}
x&=(1-\rho)a_0+\rho y_1,\\
x&=(1-\rho)c_0+(1-\rho)\rho c_1+\rho^2y_2,\\
x&=(1-\rho)d_0+(1-\rho)\rho d_1
  +(1-\rho)\rho^2d_2+\rho^3y_3,\\
&\ \vdots
\end{align*}
Each digit is either $0$ or $1$, and every $y_n$ belongs to $[0,1]$.
Different letters emphasize that representations using different numbers
of digits need not initially agree in their first, second, or later digits.
We now select digits that form a single infinite sequence.

First, each representation begins with either $0$ or $1$. At least one
of these two values must occur as the first digit in infinitely many
representations. Call that value $b_0$ and keep only those representations.

Next, discard any representation with fewer than two digits. Infinitely
many remain, and their second digit is again either $0$ or $1$. Choose
a value occurring in that position infinitely often, call it $b_1$,
and keep only those representations. They now all begin with $b_0,b_1$.

Continue in the same way. At each step we discard the finitely many
representations that are too short, then choose the next digit so that
infinitely many representations agree with all digits chosen so far.
This constructs a sequence $b_0,b_1,b_2,\ldots$.

To check that its series equals $x$, fix any $m\geq1$. There are still
representations using $n\geq m$ digits whose first $m$ digits are
$b_0,\ldots,b_{m-1}$. In any such representation, everything after
these first $m$ terms, including $\rho^n y_n$, is nonnegative and at most
\begin{equation*}
(1-\rho)\sum_{j=m}^{n-1}\rho^j+\rho^n=\rho^m.
\end{equation*}
Therefore the partial sums of the single sequence we have chosen satisfy
\begin{equation*}
0\leq x-(1-\rho)\sum_{j=0}^{m-1}\rho^j b_j\leq\rho^m.
\end{equation*}
Since $\rho^m\to0$, these partial sums converge to $x$. Hence
$x=(1-\rho)\sum_{j=0}^{\infty}\rho^j b_j$, so $x\in C_\rho$.
This proves the equality.

This intersection is closed and contained in $[0,1]$, hence compact by
the Heine--Borel theorem. Finally,
$\mathcal K_\rho=\{(x,1-x):x\in C_\rho\}$ is the image of this compact
set under the continuous map $x\mapsto(x,1-x)$, so it is compact.
\end{proof}

\paragraph{Strict advantage and persistent use.}
Define the maximum advantage of product~2 on the boundary set by
\begin{equation}
\mathcal A_\rho=\max_{z\in\mathcal K_\rho}
\left\{-\Delta c+\gamma\Delta p
\left[V_\rho(T_\rho^s z)-V_\rho(T_\rho^f z)\right]\right\}.
\end{equation}
This maximum exists by continuity and compactness.

\begin{proposition}[Persistent use of product~2]
\label{prop:persistence}
If $\mathcal A_\rho>0$, then there is a constant
$\eta>0$, such that
\begin{equation}
\liminf_{N\to\infty}\frac1N\sum_{n=1}^{N}\Ind{i_n=2}\geq\eta
\qquad\text{a.s},
\end{equation}
where $i_n$ is the product chosen at the $n$th interaction with the provider~$A$.
In particular, product~2 is used infinitely often.
\end{proposition}

\begin{proof}
Fix an optimal policy and an initial state $x_0\in\mathcal Z$.
Throughout the proof, $x=(x_s,x_f)$ denotes the full normalized
state in
\begin{equation*}
\mathcal Z=\{(x_s,x_f):x_s\ge0,\ x_f\ge0,\ x_s+x_f\le1\},
\end{equation*}
and $\|\cdot\|$ denotes the Euclidean norm.
Let $x_n$ be the state after the $n$th interaction with provider A.
Thus, product $i_n$ is chosen at state $x_{n-1}$.

Since $D_\rho$ is continuous and strictly positive on the compact
set $\mathcal Z$, it is bounded below by a positive constant.
Consequently, each successive interaction with provider A occurs
in finite calendar time almost surely, so the interaction sequence
is well defined.

To locate states where product 2 is uniquely optimal, define,
for every $x\in\mathcal Z$,
\begin{equation*}
a_\rho(x)
:=-\Delta c+\gamma\Delta p
\left[
V_\rho(T_\rho^s(x))-V_\rho(T_\rho^f(x))
\right].
\end{equation*}
By continuity of $V_\rho$ and the transition maps,
$a_\rho$ is continuous on $\mathcal Z$.
Since $\mathcal K_\rho$ is compact, the maximum
\begin{equation*}
\mathcal A_\rho
=\max_{x\in\mathcal K_\rho}a_\rho(x)
\end{equation*}
is attained. Choose $x^*\in\mathcal K_\rho$ such that
\begin{equation*}
a_\rho(x^*)=\mathcal A_\rho>0.
\end{equation*}

By continuity at $x^*$, there exists $\varepsilon>0$ such that,
for all $x\in\mathcal Z$,
\begin{equation*}
\|x-x^*\|<\varepsilon
\quad\Longrightarrow\quad
|a_\rho(x)-\mathcal A_\rho|
<\frac{\mathcal A_\rho}{2}.
\end{equation*}
In particular,
\begin{equation}
\label{eq:pers-neighborhood}
\|x-x^*\|<\varepsilon
\quad\Longrightarrow\quad
a_\rho(x)>
\mathcal A_\rho-\frac{\mathcal A_\rho}{2}
=\frac{\mathcal A_\rho}{2}>0.
\end{equation}

Here $a_\rho(x)$ is the difference between the values of choosing
products 2 and 1, conditional on a purchase from provider A and
with optimal continuation in both cases.
Therefore, product 2 is uniquely optimal at every state in this
neighborhood.

We next show that a single finite outcome sequence brings every
starting state into this neighborhood. For any
$x,\widetilde x\in\mathcal Z$, applying the success transition
to both states gives
\begin{equation*}
\begin{aligned}
T_\rho^s(x)-T_\rho^s(\widetilde x)
&=(\rho x_s+1-\rho,\rho x_f)
 -(\rho\widetilde x_s+1-\rho,\rho\widetilde x_f)\\
&=\rho(x_s-\widetilde x_s,x_f-\widetilde x_f)\\
&=\rho(x-\widetilde x).
\end{aligned}
\end{equation*}
The same identity holds for a failure.
Hence, for either $o\in\{s,f\}$,
\begin{equation*}
\|T_\rho^o(x)-T_\rho^o(\widetilde x)\|
=\rho\|x-\widetilde x\|.
\end{equation*}

For a sequence $w=(w_1,\ldots,w_m)\in\{s,f\}^m$,
listed in chronological order, write
\begin{equation*}
T_\rho^w
:=T_\rho^{w_m}\circ\cdots\circ T_\rho^{w_1}.
\end{equation*}
Repeated application of the preceding identity yields
\begin{equation}
\label{eq:pers-contraction}
\|T_\rho^w(x)-T_\rho^w(\widetilde x)\|
=\rho^m\|x-\widetilde x\|.
\end{equation}

Since the diameter of $\mathcal Z$ is $\sqrt{2}$,
choose an integer $m\ge1$ such that
\begin{equation*}
\sqrt{2}\,\rho^m<\varepsilon.
\end{equation*}
This choice of $m$ does not depend on the starting state.

The invariance identity
\begin{equation*}
\mathcal K_\rho
=T_\rho^s(\mathcal K_\rho)\cup T_\rho^f(\mathcal K_\rho),
\end{equation*}
iterated $m$ times, gives
\begin{equation*}
\mathcal K_\rho
=\bigcup_{w\in\{s,f\}^m}T_\rho^w(\mathcal K_\rho).
\end{equation*}
Because $x^*\in\mathcal K_\rho$, there exist a sequence
$w\in\{s,f\}^m$ and a state
$\widetilde x\in\mathcal K_\rho$ such that
\begin{equation*}
T_\rho^w(\widetilde x)=x^*.
\end{equation*}

Fix this sequence $w$.
For every starting state $x\in\mathcal Z$,
\begin{equation}
\label{eq:pers-word}
\begin{aligned}
\|T_\rho^w(x)-x^*\|
&=\|T_\rho^w(x)-T_\rho^w(\widetilde x)\|\\
&=\rho^m\|x-\widetilde x\|\\
&\le\sqrt{2}\,\rho^m<\varepsilon.
\end{aligned}
\end{equation}
Thus, the same sequence $w$ brings every starting state into
the neighborhood in \eqref{eq:pers-neighborhood}.

To bound the probability of this sequence conditional on past
observations, let $\mathcal F_n$ denote the complete history
through the outcome of interaction $n$, including the initial
state and past product choices.
Set
\begin{equation*}
q:=\min\{p_1,1-p_2\}>0.
\end{equation*}

At every interaction, regardless of the chosen product,
the conditional probability of success is at least $p_1$,
and that of failure is at least $1-p_2$.
Therefore,
\begin{equation*}
\begin{aligned}
\Pr(\text{success at interaction }n+1\mid\mathcal F_n)
&\ge p_1\ge q,\\
\Pr(\text{failure at interaction }n+1\mid\mathcal F_n)
&\ge1-p_2\ge q,
\end{aligned}
\qquad\text{almost surely}.
\end{equation*}

By successive conditioning, the conditional probability that
the next $m$ outcomes equal the fixed sequence $w$ is at least
$q^m$, regardless of the previous history.
This does not require independence between outcomes or blocks.

Divide the interactions into consecutive, nonoverlapping blocks
of length $m$.
Block $k$ contains interactions
\begin{equation*}
(k-1)m+1,\ldots,km.
\end{equation*}
Define the indicator
\begin{equation*}
X_k
:=\mathbf 1\{
\text{the outcomes in block }k\text{ equal }w
\}.
\end{equation*}
Then
\begin{equation}
\label{eq:pers-block-bound}
\mathbb E[X_k\mid\mathcal F_{(k-1)m}]
=\Pr(X_k=1\mid\mathcal F_{(k-1)m})
\ge q^m
\qquad\text{almost surely}.
\end{equation}

To turn \eqref{eq:pers-block-bound} into an almost-sure frequency
bound, define $M_0=0$ and, for $k\ge1$,
\begin{equation*}
M_k
:=\sum_{j=1}^{k}
\left(
X_j-\mathbb E[X_j\mid\mathcal F_{(j-1)m}]
\right).
\end{equation*}
We verify that $(M_k)_{k\ge0}$ is a martingale with respect to
$(\mathcal F_{km})_{k\ge0}$.

First, $X_j$ is $\mathcal F_{jm}$-measurable and
$\mathbb E[X_j\mid\mathcal F_{(j-1)m}]$ is
$\mathcal F_{(j-1)m}$-measurable.
For $j\le k$,
\begin{equation*}
\mathcal F_{(j-1)m}
\subseteq\mathcal F_{jm}
\subseteq\mathcal F_{km}.
\end{equation*}
Thus, every term in $M_k$ is determined by the history available
after block $k$, so $M_k$ is $\mathcal F_{km}$-measurable.

Second, since $X_j$ and its conditional expectation belong
to $[0,1]$,
\begin{equation*}
\left|
X_j-\mathbb E[X_j\mid\mathcal F_{(j-1)m}]
\right|\le1.
\end{equation*}
Consequently,
\begin{equation*}
\mathbb E[|M_k|]
\le
\mathbb E\!\left[
\sum_{j=1}^{k}
\left|
X_j-\mathbb E[X_j\mid\mathcal F_{(j-1)m}]
\right|
\right]
\le k<\infty.
\end{equation*}
In fact, $|M_k|\le k$, so
$\mathbb E[M_k^2]\le k^2<\infty$ as well.

Third, separating the last increment and taking conditional
expectations gives
\begin{equation*}
\begin{aligned}
\mathbb E[M_k\mid\mathcal F_{(k-1)m}]
&=M_{k-1}
+\mathbb E\!\left[
X_k-\mathbb E[X_k\mid\mathcal F_{(k-1)m}]
\mid\mathcal F_{(k-1)m}
\right]\\
&=M_{k-1}
+\mathbb E[X_k\mid\mathcal F_{(k-1)m}]
-\mathbb E[X_k\mid\mathcal F_{(k-1)m}]\\
&=M_{k-1}.
\end{aligned}
\end{equation*}
This establishes the martingale property.

Also, the tower property gives
\begin{equation*}
\mathbb E\!\left[
X_j-\mathbb E[X_j\mid\mathcal F_{(j-1)m}]
\right]
=\mathbb E[X_j]-\mathbb E[X_j]=0,
\end{equation*}
so $\mathbb E[M_k]=0$ for every $k$.

To obtain almost-sure convergence, observe that
\begin{equation*}
|M_k-M_{k-1}|
=
\left|
X_k-\mathbb E[X_k\mid\mathcal F_{(k-1)m}]
\right|
\le1.
\end{equation*}
Hence
\begin{equation*}
\mathbb E[(M_k-M_{k-1})^2]\le1,
\end{equation*}
and therefore
\begin{equation*}
\sum_{k=1}^{\infty}
\frac{\mathbb E[(M_k-M_{k-1})^2]}{k^2}
\le
\sum_{k=1}^{\infty}\frac{1}{k^2}
<\infty.
\end{equation*}
The martingale strong law, in the form of
\cite[Exercise~4.4.11]{Durrett2019} with $b_k=k$,
now yields
\begin{equation}
\label{eq:pers-martingale-limit}
\frac{M_k}{k}\longrightarrow0
\qquad\text{almost surely}.
\end{equation}

By the definition of $M_k$ and \eqref{eq:pers-block-bound},
\begin{equation*}
\begin{aligned}
\frac{1}{k}\sum_{j=1}^{k}X_j
&=
\frac{1}{k}\sum_{j=1}^{k}
\mathbb E[X_j\mid\mathcal F_{(j-1)m}]
+\frac{M_k}{k}\\
&\ge q^m+\frac{M_k}{k}.
\end{aligned}
\end{equation*}
Taking lower limits and using
\eqref{eq:pers-martingale-limit}, we obtain
\begin{equation}
\label{eq:pers-block-frequency}
\liminf_{k\to\infty}
\frac{1}{k}\sum_{j=1}^{k}X_j
\ge q^m
\qquad\text{almost surely}.
\end{equation}

It remains to relate these block occurrences to choices of product 2.
If $X_j=1$, then \eqref{eq:pers-word} implies
\begin{equation*}
\|x_{jm}-x^*\|<\varepsilon.
\end{equation*}
By \eqref{eq:pers-neighborhood}, product 2 is uniquely optimal
at $x_{jm}$, and the optimal policy must therefore choose it
at interaction $jm+1$.
Thus,
\begin{equation*}
\mathbf 1\{i_{jm+1}=2\}\ge X_j.
\end{equation*}

For a horizon of $N$ interactions, this guaranteed product-2
choice is included in the count only when
\begin{equation*}
jm+1\le N
\quad\Longleftrightarrow\quad
j\le\frac{N-1}{m}.
\end{equation*}
Define
\begin{equation*}
K_N:=\left\lfloor\frac{N-1}{m}\right\rfloor.
\end{equation*}
The interactions $jm+1$, for $1\le j\le K_N$, are distinct
and belong to the first $N$ interactions.
Consequently,
\begin{equation*}
\sum_{n=1}^{N}\mathbf 1\{i_n=2\}
\ge
\sum_{j=1}^{K_N}\mathbf 1\{i_{jm+1}=2\}
\ge
\sum_{j=1}^{K_N}X_j.
\end{equation*}

For $N\ge m+1$, so that $K_N\ge1$, dividing by $N$ gives
\begin{equation*}
\frac{1}{N}\sum_{n=1}^{N}\mathbf 1\{i_n=2\}
\ge
\frac{K_N}{N}
\left(
\frac{1}{K_N}\sum_{j=1}^{K_N}X_j
\right).
\end{equation*}
Since
\begin{equation*}
\frac{N-1}{mN}-\frac{1}{N}
<
\frac{K_N}{N}
\le
\frac{N-1}{mN},
\end{equation*}
we have $K_N/N\to1/m$ and $K_N\to\infty$.

The averages
$K_N^{-1}\sum_{j=1}^{K_N}X_j$ belong to $[0,1]$,
so taking lower limits yields
\begin{equation*}
\begin{aligned}
\liminf_{N\to\infty}
\frac{1}{N}\sum_{n=1}^{N}\mathbf 1\{i_n=2\}
&\ge
\liminf_{N\to\infty}
\frac{K_N}{N}
\left(
\frac{1}{K_N}\sum_{j=1}^{K_N}X_j
\right)\\
&=
\frac{1}{m}
\liminf_{N\to\infty}
\frac{1}{K_N}\sum_{j=1}^{K_N}X_j\\
&\ge
\frac{q^m}{m}>0
\qquad\text{almost surely},
\end{aligned}
\end{equation*}
where the final inequality follows from
\eqref{eq:pers-block-frequency}.

Taking $\eta=q^m/m$ proves the assertion.
Since finitely many uses of product 2 would give a limiting
frequency of zero, product 2 is used infinitely often
almost surely.
\end{proof}

\begin{proposition}[Characterization of positive advantage]
\label{prop:positive-characterization}
Let $U^{1}_\rho$ be the value of always choosing product~1.
For $\rho\in(0,1)$, we have $\mathcal A_\rho>0$ if and only if
\begin{equation}
\max_{z\in\mathcal K_\rho}
\left[U_\rho^1(T_\rho^s z)-U_\rho^1(T_\rho^f z)\right]
>\frac{\Delta c}{\gamma\Delta p}.
\end{equation}
\end{proposition}

\begin{proof}
Conditional on a purchase from seller~$A$, let
\begin{equation*}
Q_\rho^i(z)
=R-c_i+\gamma\left[p_iU_\rho^1(T_\rho^s z)
+(1-p_i)U_\rho^1(T_\rho^f z)\right],
\qquad i\in\{1,2\}.
\end{equation*}
This is the value of choosing product~$i$ for the current purchase and
product~1 thereafter. Subtracting the two expressions gives
\begin{equation*}
Q_\rho^2(z)-Q_\rho^1(z)
=-\Delta c+\gamma\Delta p
\left[U_\rho^1(T_\rho^s z)-U_\rho^1(T_\rho^f z)\right].
\end{equation*}
Define
\begin{equation*}
\mathcal A_\rho^1
=\max_{z\in\mathcal K_\rho}
\left\{-\Delta c+\gamma\Delta p
\left[U_\rho^1(T_\rho^s z)-U_\rho^1(T_\rho^f z)\right]\right\}.
\end{equation*}
The fixed-policy operator is a contraction preserving continuity,
so $U_\rho^1$ is continuous and this maximum exists.
We first show that $\mathcal A_\rho>0$ if and only if
$\mathcal A_\rho^1>0$, by proving that their nonpositivity conditions
are equivalent.

Suppose $\mathcal A_\rho^1\leq0$. Then
$ Q_\rho^1(z)\geq Q_\rho^2(z)$ at every
$z\in\mathcal K_\rho$. By the fixed-policy Bellman equation,
\begin{equation*}
U_\rho^1(z)
=\frac{D_\rho(z)}{1-\gamma[1-D_\rho(z)]}Q_\rho^1(z)
=\frac{D_\rho(z)}{1-\gamma[1-D_\rho(z)]}
\max_{i\in\{1,2\}} Q_\rho^i(z).
\end{equation*}
Thus $U_\rho^1$ also satisfies the optimal Bellman equation on
$\mathcal K_\rho$. Since this set is invariant under both transitions,
the Bellman operator restricted to it is a contraction. Uniqueness
therefore gives $U_\rho^1=V_\rho$ on $\mathcal K_\rho$, and hence
$\mathcal A_\rho=\mathcal A_\rho^1\leq0$.

Conversely, suppose $\mathcal A_\rho\leq0$. Product~1 then maximizes
the Bellman expression evaluated at $V_\rho$ at every state in
$\mathcal K_\rho$. Consequently, $V_\rho$ satisfies the Bellman equation
for always choosing product~1 on this invariant set. Uniqueness of
the fixed-policy value gives $V_\rho=U_\rho^1$ there, and hence
$\mathcal A_\rho^1=\mathcal A_\rho\leq0$.

We have proved $\mathcal A_\rho\leq0$ if and only if
$\mathcal A_\rho^1\leq0$. Taking negations gives the desired equivalence
of strict positivity. Finally, since $\Delta c$ and $\gamma\Delta p>0$
do not depend on $z$,
\begin{equation*}
\mathcal A_\rho^1
=-\Delta c+\gamma\Delta p\max_{z\in\mathcal K_\rho}
\left[U_\rho^1(T_\rho^s z)-U_\rho^1(T_\rho^f z)\right].
\end{equation*}
Therefore $\mathcal A_\rho^1>0$ is equivalent to the inequality
in the statement.
\end{proof}

% Let $V_\rho$ be the unique bounded continuous optimal value function on
% $\mathcal Z$, whose existence and continuity are established by the
% manuscript's contraction argument. Define its continuation gap by
% \[
% G_\rho(z)=\gamma\left[V_\rho(T_\rho^s z)-V_\rho(T_\rho^f z)\right].
% \]
% Conditional on provider~$A$ being selected, product~2 is strictly optimal at
% $z$ precisely when
% \begin{equation}\label{eq:strict}
% -\Delta c+\Delta p\,G_\rho(z)>0.
% \end{equation}
% An optimal policy below means a measurable selection of a Bellman-maximizing
% product at every state. Either maximizing product may be selected at ties.

% Let $\mathcal K_\rho$ be the manuscript's limiting reachable set. Equivalently,
% it is the nonempty compact attractor of the two outcome maps:
% \begin{equation}\label{eq:attractor}
% \mathcal K_\rho=T_\rho^s(\mathcal K_\rho)
% \mathbin{\cup}T_\rho^f(\mathcal K_\rho).
% \end{equation}
% Both maps contract sup-norm distances by the factor $\rho$. The attractor is
% contained in $x_s+x_f=1$; it is the full boundary segment when
% $\rho\geq1/2$, and need not be an interval when $\rho<1/2$.
% These compactness and invariance claims are established in
% Lemma~\ref{lem:compactness} below.

% \begin{lemma}[Compactness and the range of forgetting parameters]
% \label{lem:compactness}
% For every fixed $\rho\in[0,1)$, $\mathcal Z$ is a compact invariant state
% space, and $\mathcal K_\rho$ is a nonempty compact invariant subset of the
% boundary $x_s+x_f=1$. In original coordinates the closed domain
% \[
% \left\{(S,F):S\geq0,\ F\geq0,\ S+F\leq\frac1{1-\rho}\right\}
% \]
% is also compact. These statements do not require $\rho\geq1/2$.
% \end{lemma}

% \begin{proof}
% The set $\mathcal Z$ is the intersection of three closed half-spaces and is
% contained in $[0,1]^2$. It is therefore closed and bounded, hence compact
% by the Heine--Borel theorem. For either outcome, the coordinates remain
% nonnegative and their sum becomes
% \[
% \rho(x_s+x_f)+1-\rho\leq1.
% \]
% Thus both continuous outcome maps send $\mathcal Z$ into itself.

% We parameterize the boundary $x_s+x_f=1$ by $x=x_s\in[0,1]$, so that
% $x_f=1-x$. On this boundary, define the two maps
% \begin{equation}
% f_0(x)=\rho x,
% \qquad
% f_1(x)=\rho x+(1-\rho).
% \end{equation}
% Starting from
% \begin{equation}
% K_0=[0,1],
% \end{equation}
% define recursively
% \begin{equation}
% K_{n+1}=f_0(K_n)\cup f_1(K_n).
% \end{equation}
% Write
% \begin{equation}
% C_\rho=\bigcap_{n=0}^{\infty}K_n.
% \end{equation}
% We will show that $C_\rho$ is nonempty and compact, and identify
% $\mathcal K_\rho=\{(x,1-x):x\in C_\rho\}$.

% We consider separately the cases $\rho\geq1/2$ and $\rho<1/2$.

% \medskip
% \noindent\textit{Case 1: $\rho\geq1/2$.}

% The images of the initial interval under the two maps are
% \begin{equation}
% f_0([0,1])=[0,\rho],
% \qquad
% f_1([0,1])=[1-\rho,1].
% \end{equation}
% Since $1-\rho\leq\rho$, the two intervals overlap or meet at a single
% point, and together they cover the entire interval $[0,1]$. Therefore,
% \begin{equation}
% K_1=f_0(K_0)\cup f_1(K_0)
% =[0,\rho]\cup[1-\rho,1]=[0,1]=K_0.
% \end{equation}
% Applying the same maps again produces the same set, so
% \begin{equation}
% K_n=[0,1]\qquad\text{for every }n\geq0.
% \end{equation}
% It follows that
% \begin{equation}
% C_\rho=\bigcap_{n=0}^{\infty}K_n=[0,1],
% \end{equation}
% which is nonempty and compact.

% \medskip
% \noindent\textit{Case 2: $\rho<1/2$.}

% In this case $\rho<1-\rho$, so the two images are separated:
% \begin{equation}
% K_1=[0,\rho]\cup[1-\rho,1].
% \end{equation}
% We first show that every $K_n$ is compact. The initial set $K_0=[0,1]$
% is compact. If $K_n$ is compact, continuity of $f_0$ and $f_1$ implies
% that $f_0(K_n)$ and $f_1(K_n)$ are compact. Their finite union
% $K_{n+1}$ is therefore compact. By induction, every $K_n$ is compact.

% We next show that the sets are nested. Since both maps send $[0,1]$
% into itself,
% \begin{equation}
% K_1\subseteq K_0.
% \end{equation}
% If $K_n\subseteq K_{n-1}$, then
% \begin{equation}
% f_0(K_n)\subseteq f_0(K_{n-1}),
% \qquad
% f_1(K_n)\subseteq f_1(K_{n-1}),
% \end{equation}
% and therefore
% \begin{equation}
% K_{n+1}=f_0(K_n)\cup f_1(K_n)
% \subseteq f_0(K_{n-1})\cup f_1(K_{n-1})=K_n.
% \end{equation}
% Hence,
% \begin{equation}
% K_0\supseteq K_1\supseteq K_2\supseteq\cdots.
% \end{equation}
% Each $K_n$ is closed, so their intersection $C_\rho$ is closed.
% Moreover, $C_\rho\subseteq[0,1]$, so it is bounded. By the
% Heine--Borel theorem, $C_\rho$ is compact. It is nonempty because
% $0$ and $1$ belong to every $K_n$.

% \medskip
% In both cases, the continuous image
% \begin{equation}
% \{(x,1-x):x\in C_\rho\}
% \end{equation}
% is nonempty and compact. To check invariance, nesting gives
% $f_0(C_\rho)\cup f_1(C_\rho)\subseteq C_\rho$.
% Conversely, for $x\in C_\rho$, write $x=f_{a_n}(y_n)$ with
% $a_n\in\{0,1\}$ and $y_n\in K_n$. Choose a subsequence with constant
% $a_n$ and a further subsequence with $y_n\to y$. Nesting and closedness
% imply $y\in C_\rho$, and continuity gives $x=f_{a_n}(y)$.
% Thus $C_\rho=f_0(C_\rho)\cup f_1(C_\rho)$.

% This set also gives precisely the limiting reachable states.
% Each $K_n$ is the union of images of $[0,1]$ under all length-$n$
% outcome words. For $x\in C_\rho$, choose such a word whose image
% contains $x$. The corresponding two-dimensional word maps a boundary
% point to $(x,1-x)$. Applying it instead to $(0,0)$ gives a reachable
% state within $\rho^n$ in the supremum norm. Conversely, a state reached
% from $(0,0)$ after $m$ observations is within $\rho^m$ of the same word
% applied to a boundary point. For $m\geq n$, that boundary image has its
% first coordinate in $K_m\subseteq K_n$. Thus any limit as history lengths
% tend to infinity belongs to $\{(x,1-x):x\in K_n\}$ for every $n$.
% Consequently,
% \begin{equation}
% \mathcal K_\rho=\{(x,1-x):x\in C_\rho\}.
% \end{equation}

% Finally, for fixed $\rho<1$, dividing both coordinates by $1-\rho$
% is a continuous invertible linear change of coordinates. Its image of
% $\mathcal Z$ is the original closed domain in the statement, which is
% therefore compact as well.
% \end{proof}

% \paragraph{Endpoints and finite histories.}
% At $\rho=0$, the maps are constant and
% $\mathcal K_0=\{(1,0),(0,1)\}$, so compactness still holds.
% For $0<\rho<1$, states reached from $(0,0)$ after finitely many observations
% satisfy $x_s+x_f=1-\rho^k<1$. That finite-history reachable set need not be
% closed; the compact sets above include the limiting states. No finite-time
% arrival at the boundary is asserted.
% At $\rho=1$, the normalization is not invertible and the original process
% has $S+F=k$, with no bounded state domain over all histories. Thus this
% compactness argument does not apply to cumulative learning. Also, the
% original closed domains are not uniformly bounded as $\rho\uparrow1$.
% The fixed geometric triangle $\mathcal Z$ remains compact as a set, but its
% interpretation as an equivalent normalized state space requires $\rho<1$.

% Denote the maximum conditional gain on that set by
% \begin{equation}\label{eq:maxgain}
% \mathcal A_\rho
% =\max_{z\in\mathcal K_\rho}
% \left\{-\Delta c+\Delta p\,G_\rho(z)\right\}.
% \end{equation}
% The maximum exists by continuity and compactness. This is the scalar denoted
% $M(\rho)$ in the companion analysis note; it is not the manuscript's
% state-dependent conditional value $M_\rho(z)$.

% Let $z_k$ be the state after the $k$th purchase from~$A$, and let $i_k$ be
% the product supplied at that purchase. Thus $i_k$ is chosen at $z_{k-1}$.
% Conditional on the history and $i_k=i$, the outcome is a success with
% probability $p_i$, and $z_k$ is respectively $T_\rho^s z_{k-1}$ or
% $T_\rho^f z_{k-1}$. These conditional probabilities, rather than independence
% of realized product choices, are what the proof uses.

% \begin{proposition}[Persistent high-quality product use]\label{prop:persistence}
% Suppose $\mathcal A_\rho>0$. There exists $\eta>0$, depending on the model
% parameters but not on the initial state, such that under any optimal policy
% as specified above, for every $z_0\in\mathcal Z$,
% \begin{equation}\label{eq:persistence}
% \liminf_{N\to\infty}\frac1N\sum_{k=1}^N\Ind{i_k=2}
% \geq\eta
% \qquad\text{almost surely}.
% \end{equation}
% In particular, the conclusion holds from $z_0=(0,0)$, without requiring the
% process to reach the boundary in finite time.
% \end{proposition}

% \begin{proof}
% \textit{Step 1: a finite outcome sequence forces a strictly optimal
% product-2 decision.}
% By \eqref{eq:maxgain}, there is $z^*\in\mathcal K_\rho$ satisfying
% \eqref{eq:strict}. Continuity of $G_\rho$ gives $\delta>0$ such that
% \begin{equation}\label{eq:neighborhood}
% z\in\mathcal Z,\quad \|z-z^*\|_\infty<\delta
% \quad\Longrightarrow\quad
% -\Delta c+\Delta p\,G_\rho(z)>0.
% \end{equation}
% Choose an integer $L\geq1$ with $\rho^L<\delta$.
% Iterating \eqref{eq:attractor}, the images of $\mathcal K_\rho$ under all
% outcome words of length~$L$ cover $\mathcal K_\rho$. Hence there are
% $w_1,\ldots,w_L\in\{s,f\}$ and $y\in\mathcal K_\rho$ such that
% \[
% z^*=T_\rho^{w_L}\circ\cdots\circ T_\rho^{w_1}(y).
% \]
% For every $z\in\mathcal Z$, contraction and
% $\operatorname{diam}_\infty(\mathcal Z)=1$ give
% \begin{equation}\label{eq:forcing}
% \left\|T_\rho^{w_L}\circ\cdots\circ T_\rho^{w_1}(z)-z^*\right\|_\infty
% =\rho^L\|z-y\|_\infty\leq\rho^L<\delta.
% \end{equation}
% Consequently, observing the prescribed word forces the next product choice
% to be~2 under every optimal policy, regardless of the starting state.

% \medskip
% \noindent\textit{Step 2: every block has a uniformly positive conditional
% probability of containing product~2.}
% Put
% \[
% \varepsilon=\min\{p_1,1-p_2\}>0,
% \qquad m=L+1,\qquad q=\varepsilon^L>0.
% \]
% Consider any block of $m$ purchases and condition on its entire starting
% history. Whatever product is chosen, the conditional probability of success
% is at least $p_1$, and that of failure is at least $1-p_2$. Thus every
% prescribed outcome has conditional probability at least $\varepsilon$,
% including after any prefix of the word $w_1,\ldots,w_L$.
% Repeated conditioning therefore gives
% \[
% \Pb(\text{the first $L$ outcomes of the block are }w_1,\ldots,w_L
% \mid\text{history before the block})\geq\varepsilon^L=q.
% \]
% Whenever this word occurs, \eqref{eq:forcing} and \eqref{eq:neighborhood}
% force product~2 at the last, $(L+1)$st, purchase of the block. Consequently,
% \begin{equation}\label{eq:blockbound}
% \Pb(\text{the block contains product~2}\mid
% \text{history before the block})\geq q.
% \end{equation}
% The strict assumption $p_2<1$ ensures that $\varepsilon>0$. No independence
% of product choices or outcome blocks is needed.

% \medskip
% \noindent\textit{Step 3: the block bound implies a positive lower frequency.}
% Partition the purchases into disjoint blocks of length $m$.
% Let $Y_j$ indicate that block~$j$ contains at least one product-2 choice,
% and let $\mathcal F_j$ be the history through purchase $jm$.
% By \eqref{eq:blockbound},
% \[
% \E[Y_j\mid\mathcal F_{j-1}]\geq q.
% \]
% The variables $Y_j-\E[Y_j\mid\mathcal F_{j-1}]$ are martingale
% differences bounded in absolute value by~1. For every $a>0$, the
% Azuma--Hoeffding inequality gives
% \[
% \Pb\left(
% \left|\sum_{j=1}^J
% \bigl(Y_j-\E[Y_j\mid\mathcal F_{j-1}]\bigr)\right|\geq aJ
% \right)\leq2\exp(-a^2J/2).
% \]
% The right-hand side is summable in $J$. Borel--Cantelli, applied for
% $a=1,1/2,1/3,\ldots$, therefore yields
% \[
% \frac1J\sum_{j=1}^J
% \bigl(Y_j-\E[Y_j\mid\mathcal F_{j-1}]\bigr)\longrightarrow0
% \quad\text{almost surely}.
% \]
% It follows that
% \[
% \liminf_{J\to\infty}\frac1J\sum_{j=1}^J Y_j\geq q
% \quad\text{almost surely}.
% \]
% For arbitrary $N$, set $J=\lfloor N/m\rfloor$. Each successful complete
% block contributes at least one distinct product-2 purchase, so
% \[
% \frac1N\sum_{k=1}^N\Ind{i_k=2}
% \geq\frac{J}{N}\,\frac1J\sum_{j=1}^JY_j.
% \]
% Since $J/N\to1/m$, \eqref{eq:persistence} follows with
% \[
% \eta=\frac{q}{m}=
% \frac{\bigl(\min\{p_1,1-p_2\}\bigr)^L}{L+1}>0.
% \]
% \end{proof}

% \paragraph{Interpretation and scope.}
% The result concerns the exact policy on the full state space, not a numerical
% policy obtained by projecting the state onto the boundary. The neighborhood
% in \eqref{eq:neighborhood} includes interior states, and \eqref{eq:forcing}
% holds for every state in $\mathcal Z$. Thus the proof does not assume finite-time
% arrival at the boundary. Nor does it require a single-threshold policy,
% independence between blocks, or a unique invariant distribution.

% The positive lower bound may be extremely small. The proposition establishes
% neither convergence of the empirical frequency to a unique constant nor a
% useful finite-horizon visitation rate. Identifying parameter values for which
% $\mathcal A_\rho>0$ is a separate task; that inequality is the hypothesis,
% not a conclusion asserted for all intermediate forgetting rates.

% Finally, purchases from~$A$ occur infinitely often almost surely in the
% calendar-time model. For fixed $\rho<1$, its Beta-tail demand $D_\rho$ is
% continuous and strictly positive on compact $\mathcal Z$, hence bounded below
% by a positive constant. Conditional on any history, the probability of no
% purchase from~$A$ in the next $n$ calendar periods is bounded by a geometric
% term tending to zero. The interaction index used above is therefore
% well-defined for all $k$ almost surely. Statement~\eqref{eq:persistence}
% measures the frequency among these purchases.

\begin{thebibliography}{1}
\bibitem{Durrett2019}
Rick Durrett.
\textit{Probability: Theory and Examples}.
Fifth edition. Cambridge University Press, 2019.
\end{thebibliography}

\end{document}
